\documentclass[a4paper]{article}
% Español (México) con XeLaTeX: fontspec para fuentes Unicode, polyglossia
% para la división silábica y los nombres del documento en la variante mexicana.
\usepackage{fontspec}
\usepackage{polyglossia}
\setdefaultlanguage[variant=mexican]{spanish}
% Los nombres chinos van como «pinyin (original)»: xeCJK da a los caracteres
% Han su propia fuente; sin ella XeLaTeX los omite con un aviso.
% FandolSong no cubre algunos caracteres de nombres (U+73FA); WenQuanYi Zen Hei sí.
\usepackage[AutoFallBack=true]{xeCJK}
\setCJKmainfont{FandolSong-Regular.otf}
\setCJKfallbackfamilyfont{\CJKrmdefault}{WenQuanYi Zen Hei}
% polyglossia no traduce los nombres de listings.
\AtBeginDocument{\providecommand{\lstlistingname}{}\renewcommand{\lstlistingname}{Listado}%
  \providecommand{\lstlistlistingname}{}\renewcommand{\lstlistlistingname}{Índice de listados}}
\usepackage{amsmath, amssymb}
\usepackage{graphicx}
\usepackage[margin=2.35cm]{geometry}
\usepackage[most]{tcolorbox}
\usepackage{booktabs}
\usepackage{tabularx}
\usepackage{longtable}
\usepackage{array}
\usepackage{float}
\usepackage{xcolor}
\usepackage{hyperref}
\usepackage{fancyhdr}
\setCJKmainfont[AutoFakeBold=2.4]{AR PL UMing CN}
\setCJKsansfont[AutoFakeBold=2.4]{Droid Sans Fallback}
\setCJKmonofont[AutoFakeBold=2.4]{Droid Sans Fallback}
\hypersetup{colorlinks=true, linkcolor=blue!55!black, urlcolor=blue!60!black}
\graphicspath{{./}{figures/}}
\setlength{\parskip}{0.55em}
\setlength{\parindent}{2em}
\setlength{\emergencystretch}{3em}
\renewcommand{\arraystretch}{1.18}
\newtcolorbox{knowledgebox}[1]{enhanced, colback=blue!5!white, colframe=blue!70!black, colbacktitle=blue!70!black, coltitle=white, fonttitle=\bfseries, title={#1}, sharp corners, breakable}
\newtcolorbox{importantbox}[1]{enhanced, colback=yellow!10!white, colframe=yellow!75!black, colbacktitle=yellow!75!black, coltitle=black, fonttitle=\bfseries, title={#1}, sharp corners, breakable}
\newtcolorbox{warningbox}[1]{enhanced, colback=red!5!white, colframe=red!70!black, colbacktitle=red!70!black, coltitle=white, fonttitle=\bfseries, title={#1}, sharp corners, breakable}
\pagestyle{fancy}
\fancyhf{}
\fancyfoot[C]{\thepage}
\fancyfoot[R]{\footnotesize\color{gray}\href{https://github.com/hqhq1025/ai-course-notes}{AI Course Notes}}
\renewcommand{\headrulewidth}{0pt}
\renewcommand{\footrulewidth}{0pt}
\newcommand{\videourl}{https://www.youtube.com/watch?v=40qPt8R2uys}
\newcommand{\repourl}{https://github.com/hqhq1025/ai-course-notes}

\begin{document}
\begin{titlepage}
\centering
{\Large Notas didácticas de una entrevista a fondo\par}
\vspace{1.0cm}
{\huge\bfseries Language es veneno, eliminar la L y la transición de XPeng (小鹏) hacia la Physical AI\par}
\vspace{0.5cm}
{\Large Zhang Xiaojun conversa con Liu Xianming (刘先明), responsable de conducción autónoma de XPeng\par}
\vspace{0.35cm}
{\large Elaborado a partir del video público de Zhang Xiaojun Podcast, la transcripción hecha en GPU, la estructura de capítulos y diagramas conceptuales propios\par}
\vspace{0.3cm}
{\large 2025-11-18\par}
\vspace{0.85cm}
\includegraphics[width=0.88\textwidth,height=0.42\textheight,keepaspectratio]{cover.jpg}\par
\vfill
\begin{tcolorbox}[width=0.92\textwidth, colback=black!2!white, colframe=black!55, sharp corners]
\textbf{Título del video}: 120. Primera entrevista con Liu Xianming, recién nombrado en XPeng: Language es veneno, eliminar la L, lo simple es bello, el cambio de líder y la transformación de XPeng hacia la IA\par
\textbf{Canal del video}: Zhang Xiaojun Podcast\par
\textbf{Duración del video}: 01:48:46\par
\textbf{Enlace del video}: \href{\videourl}{\nolinkurl{\videourl}}\par
\textbf{Notas sobre el material}: YouTube no tiene subtítulos disponibles; el texto se basa en la transcripción con faster-whisper large-v3 en CUDA, los capítulos del video, la portada y figuras didácticas propias. Las tomas fijas de la entrevista no se reutilizan en el texto.\par
\textbf{Página del proyecto}: \href{\repourl}{\nolinkurl{\repourl}}\par
\end{tcolorbox}
\end{titlepage}

\tableofcontents
\newpage
\section{Introducción: por qué una automotriz pasa de la conducción autónoma a la Physical AI}

Esta sección expone primero la pregunta central del episodio. Liu Xianming (刘先明) asumió la dirección del centro de conducción autónoma de XPeng (小鹏), y la atención externa se concentró en el cambio de liderazgo; sin embargo, lo que realmente vale la pena documentar de la entrevista es que él vuelve a entender la conducción autónoma dentro del marco de la Physical AI: el automóvil no es solo un medio de transporte, sino también un agente inteligente en el mundo real; la conducción autónoma no es solo un sistema de percepción, planificación y control, y mapas, sino un ciclo cerrado inteligente formado en conjunto por los datos, el modelo, la fábrica en la nube, el despliegue en el vehículo y la retroalimentación de la producción en serie.

El juicio más tajante del episodio es «Language es veneno». No quiere decir que el lenguaje carezca de valor, sino que en la cadena de entrenamiento de la conducción autónoma o de la IA física, si las señales de los sensores se traducen primero a tokens de un lenguaje intermedio y luego se decodifica la trayectoria a partir del lenguaje, se introducen supervisión humana, un cuello de botella de compresión y obstáculos para el data scaling. La ruta de Liu Xianming consiste en ir eliminando capas intermedias: eliminar la dependencia del lidar, eliminar las reglas de planificación y control, eliminar la estructura intermedia del sistema de extremo a extremo y, por último, eliminar también el language bottleneck.

\begin{importantbox}{Tesis central del episodio}
La transformación de XPeng hacia la IA no consiste en «agregarle un modelo grande al automóvil», sino en reescribir la conducción autónoma como Physical AI: con modelos más grandes, más datos, menos capas intermedias humanas y un ciclo cerrado de ingeniería más sólido, para que el automóvil aprenda, se despliegue e itere de forma continua en el mundo real.
\end{importantbox}

\begin{knowledgebox}{Nota sobre la estrategia visual}
Este video es una entrevista con encuadre fijo, sin slides, pizarrón ni demostraciones de producto. En el texto principal la portada se usa solo para identificar la fuente, y todas las imágenes del texto son diagramas conceptuales propios, que explican la evolución de la pila de software de conducción autónoma, la eliminación de Language, la fábrica de modelos en la nube, el ciclo cerrado de datos del fabricante de automóviles, la reducción de niveles jerárquicos en la organización y la misión del cambio de liderazgo.
\end{knowledgebox}

\subsection{Resumen de la sección}

El eje del EP120 es la «simplificación»: en lo técnico, eliminar capas intermedias; en lo organizativo, reducir niveles jerárquicos; en lo estratégico, situar la conducción autónoma dentro de la Physical AI. No es una simple entrevista sobre un cambio de personal, sino una metodología para la transformación de una automotriz hacia la IA.
\section{La trayectoria de Liu Xianming (刘先明): de CV/ML a Robotaxi y después a un fabricante de automóviles}

Esta sección examina primero por qué Liu Xianming llegó a XPeng (小鹏). Su trayectoria incluye un doctorado en UIUC, Facebook Connectivity Lab, Horizon Robotics (地平线) en Norteamérica, Cruise Robotaxi y, por último, XPeng. Lo que da continuidad a esta línea no es un «historial de cambios de empleo», sino una elección tecnológica mission driven: de usar imágenes satelitales para atender desastres y estimar la distribución de la población mundial, a la conducción autónoma para reducir la fatiga de los conductores y los riesgos de tránsito, y después a construir, dentro de un fabricante de automóviles, un ciclo cerrado de datos controlable.

\begin{figure}[H]
\centering
\includegraphics[width=0.96\textwidth]{liuxianming-trajectory.png}
\caption{La trayectoria de Physical AI de Liu Xianming: CV/ML, Facebook, Cruise, XPeng y Physical AI. Diagrama conceptual de elaboración propia, a partir de la conversación de 00:02:16--00:19:00.}
\end{figure}

\begin{knowledgebox}{Lectura de la figura: el hilo conductor de esta trayectoria son los problemas del mundo real}
De las imágenes satelitales para atender desastres al Robotaxi, y después a la flota de vehículos de producción en serie de XPeng, Liu Xianming se ha ocupado siempre de «cómo entra la IA en el mundo real». Esto explica por qué no ve la conducción autónoma solo como una función de las empresas automotrices, sino como parte de Physical AI.
\end{knowledgebox}

\subsection{Dos lecciones de Cruise}

Dos lecciones que Liu Xianming aprendió en Cruise influyeron en su rumbo posterior. La primera es la simplificación extrema y la Infra a gran escala. Para que un Robotaxi pase de demo a producto se necesitan Data Infrastructure, Training Infrastructure, una cadena de análisis de problemas y un sistema de publicación de software. La segunda es la Continuous Learning Machine, es decir, resolver problemas de forma continua mediante la iteración con datos. Lo de «tumbarse en Hawái a esperar que caigan monedas de oro» es solo una broma; lo que hay detrás es que el ciclo cerrado de datos permite que el sistema mejore de manera continua.

\begin{importantbox}{Cómo se traslada la experiencia de Cruise}
Al pasar del Robotaxi a un fabricante de automóviles, lo decisivo no es cambiar de escenario, sino trasladar «el retorno de datos + Infra + la iteración rápida» a una flota más grande y a una cadena de datos más controlable.
\end{importantbox}

\subsection{Por qué ir a un fabricante de automóviles}

Si el futuro de la conducción autónoma y de Physical AI depende de datos a gran escala, de cadenas de datos controlables y de retroalimentación rápida, los fabricantes de automóviles tienen una ventaja natural. Las empresas de Robotaxi controlan su flota, pero su escala es limitada; los fabricantes tienen vehículos de producción en serie, usuarios reales, vialidades de distintas ciudades y una cadena completa de datos en el vehículo. Liu Xianming eligió XPeng porque él y He Xiaopeng (何小鹏) se hacían la misma pregunta: ¿cómo puede la ruta tecnológica de la siguiente generación dejar muy atrás a los competidores actuales?

\begin{knowledgebox}{Ventajas distintas de un fabricante de automóviles frente a una empresa de Robotaxi}
\begin{tabularx}{\textwidth}{p{0.22\textwidth}p{0.34\textwidth}X}
\toprule
Dimensión & Empresa de Robotaxi & Fabricante de automóviles \\
\midrule
Control de la flota & Controla la flota en operación; datos de alta calidad, pero la escala crece más despacio & Gran escala de vehículos de producción en serie; cobertura más amplia de ciudades y escenarios. \\
Objetivo comercial & Ofrece directamente servicios de movilidad & Venta de vehículos, suscripción de conducción inteligente, Robotaxi y experiencia de marca en paralelo. \\
Cadena de datos & Ciclo cerrado operativo sólido & El retorno de datos de la flota de usuarios y los problemas de la producción en serie son más complejos. \\
Restricciones de riesgo & Zona de servicio y alcance operativo controlables & Se dirige al usuario común; la responsabilidad de seguridad y calidad es mayor. \\
\bottomrule
\end{tabularx}
\end{knowledgebox}

\subsection{La presión que imponen los problemas del mundo real}

Liu Xianming comparó la presión laboral en Facebook, Cruise y XPeng. En Facebook, los problemas de investigación eran sobre todo de tecnología y recursos; en Cruise, los vehículos estaban en la calle, y los accidentes, las fallas y la seguridad eran presiones reales; en XPeng el ritmo es aún más rápido, porque un fabricante de producción en serie debe atender al mismo tiempo la IA, el producto, el negocio, la calidad, el hardware y la responsabilidad ante los usuarios. Esta presión explica por qué Physical AI no es research puro, sino que exige tomar decisiones de compromiso de forma constante dentro de sistemas reales.

\begin{warningbox}{De dónde viene la presión sobre la IA en el mundo real}
Cuando un modelo del mundo digital se equivoca, normalmente se puede revertir, reintentar u ocultar la falla; cuando un vehículo en la calle se equivoca, están en juego la seguridad, la responsabilidad, la marca y la regulación. Por eso la iteración de la IA física tiene que ser más cautelosa que la del software puro y depender más del ciclo cerrado de ingeniería.
\end{warningbox}

\subsection{De San Francisco a Cantón: el escenario cambia la estructura del problema}

Liu Xianming mencionó que, al trabajar en conducción autónoma en San Francisco, los corner case como personas en situación de calle, basura, mascotas y aves eran muy frecuentes, lo que hacía sentir que el problema era muy difícil; en cambio, en Cantón, algunos problemas que antes costaba resolver no tienen la misma magnitud. Esta observación es importante: la conducción autónoma no es un ejercicio abstracto de algoritmos; la estructura vial de la ciudad, los participantes del tránsito, la infraestructura y el comportamiento de la gente modifican la dificultad del problema.

\begin{importantbox}{El escenario no es un telón de fondo, sino parte del modelo}
El mismo algoritmo enfrenta corner case distintos en ciudades distintas. La elección del escenario afecta la distribución de los datos, el entrenamiento del modelo, el ritmo de puesta en producción y la experiencia del usuario; por eso, ni en el Robotaxi ni en la conducción inteligente de producción en serie se puede hablar solo del modelo sin hablar de la ciudad y del entorno operativo.
\end{importantbox}

\subsection{Resumen de la sección}

La trayectoria de Liu Xianming muestra que este cambio de liderazgo en XPeng no es solo un nombramiento organizacional, sino una nueva combinación de la experiencia en conducción autónoma, el ciclo cerrado de datos del Robotaxi y los recursos de producción en serie de un fabricante de automóviles.
\section{Pila de software de conducción autónoma: de las reglas a la fábrica de modelos en la nube}

La sección anterior trató sobre las personas y la organización; esta sección entra en la pila tecnológica. Liu Xianming (刘先明) divide la evolución de la conducción autónoma en varias etapas: Software 1.0 son reglas y optimización; Software 1.5 es modelos más reglas; Software 2.0 es extremo a extremo; Software 3.0, o la vía VLA/VLM, consiste en hacer scaling de datos con modelos más grandes, con más capacidad de cómputo y más datos. Al final, el hardware del vehículo no puede soportar todo el entrenamiento ni los modelos grandes, por lo que se necesita una fábrica de modelos en la nube.

\begin{figure}[H]
\centering
\includegraphics[width=0.96\textwidth]{autonomous-stack-evolution.png}
\caption{Evolución de la pila de software de conducción autónoma: de las reglas y los sistemas parcialmente basados en modelos al enfoque extremo a extremo, y de ahí a modelos más grandes y a la fábrica en la nube. Diagrama conceptual propio, elaborado a partir de la conversación de 00:02:16--00:19:00.}
\end{figure}

\begin{knowledgebox}{Lectura de la figura: cada generación derriba el límite anterior}
Software 1.0 está restringido por las reglas y la optimización; el límite de Software 1.5 sigue estando en el código de planificación y control; Software 2.0 reduce la estructura diseñada a mano mediante el enfoque extremo a extremo; Software 3.0/VLA/VLM continúa el scaling con modelos más grandes y más datos; la fábrica en la nube se encarga del entrenamiento, la destilación, la cuantización, la poda y el despliegue.
\end{knowledgebox}

\subsection{De Software 1.0 a 2.0}

Software 1.0 usa agrupamiento de datos de LiDAR, detección tradicional, reglas escritas a mano y optimización matemática para la percepción y la planificación y control. Los sistemas de esta generación son interpretables y controlables desde la ingeniería, pero quedan atrapados por el límite de las reglas. Software 1.5 introduce modelos en la percepción, por ejemplo para detección y segmentación, pero la planificación y el control siguen dependiendo de una gran cantidad de reglas. Software 2.0 busca que la red neuronal aprenda a partir de los datos una relación más completa entre entradas y salidas, y sustituye las reglas escritas a mano, cada vez más complejas, por la iteración sobre los datos.

\begin{warningbox}{Las reglas no son un error; son un límite propio de una etapa}
Los sistemas basados en reglas son muy importantes en las primeras etapas porque son controlables, ajustables e interpretables. Pero cuando la complejidad de los escenarios supera la capacidad de mantener reglas manualmente, las reglas se convierten en un límite. Desmontar las reglas no es negar la ingeniería, sino reconocer que los problemas a escala requieren un enfoque basado en datos.
\end{warningbox}

\subsection{Fábrica de modelos en la nube}

Los chips del vehículo no pueden entrenar directamente modelos de gran tamaño ni son adecuados para ejecutar directamente todas las capacidades de la nube. Por ello, Liu Xianming propone la idea de una «fábrica en la nube»: entrenar modelos enormes en la nube y luego comprimir sus capacidades mediante destilación, cuantización, poda y otros métodos, para desplegarlas en distintos tipos de hardware del vehículo. Este esquema se asemeja a una línea de producción de modelos: una vez entrenado un modelo madre en la nube, se pueden generar de forma continua modelos para el vehículo adaptados a distintas plataformas.

\begin{figure}[H]
\centering
\includegraphics[width=0.96\textwidth]{cloud-model-factory.png}
\caption{Fábrica de modelos en la nube: el modelo grande se entrena en la nube y después se destila, cuantiza y poda para llevarlo al vehículo. Diagrama conceptual propio, elaborado a partir de la conversación de 00:02:16--00:19:00.}
\end{figure}

\begin{knowledgebox}{Lectura de la figura: la nube aporta la capacidad y el vehículo, el despliegue}
El modelo grande en la nube aprovecha una gran cantidad de parámetros y grandes volúmenes de datos; después del entrenamiento, la destilación transfiere su capacidad a un modelo pequeño, que luego se cuantiza y se poda para ajustarse a la capacidad de cómputo del vehículo. Así se puede aprovechar el scaling y, al mismo tiempo, cumplir con las restricciones de hardware de los vehículos de producción en serie.
\end{knowledgebox}

\subsection{Resumen de la sección}

El núcleo de la evolución de la pila de software de conducción autónoma consiste en eliminar de forma continua las capas intermedias diseñadas a mano, ampliar la escala de los datos y de los modelos, y después usar un sistema de ingeniería para comprimir las capacidades de la nube y llevarlas al vehículo. La fábrica de modelos en la nube es la estructura fundamental que une los modelos grandes con el despliegue en producción en serie.
\section{Quitar Language: por qué Language es veneno}

Esta sección aborda el juicio técnico central del título del episodio. Muchas rutas VLA traducen primero las señales de los sensores a language tokens y después decodifican la trajectory a partir de esos language tokens. Liu Xianming (刘先明) considera que esto genera un cuello de botella, porque la representación lingüística intermedia introduce supervisión humana y compresión semántica humana, lo que perjudica el data scaling. El enfoque de Xiaopeng (小鹏) consiste en quitar esta capa intermedia de language: vision y language entran como entrada y la action se decodifica directamente.

\begin{figure}[H]
\centering
\includegraphics[width=0.94\textwidth]{remove-language-bottleneck.png}
\caption{Quitar el Language Bottleneck: los tokens lingüísticos intermedios limitan el scaling de los datos. Diagrama conceptual propio, elaborado a partir de la conversación de 00:19:00--00:33:53.}
\end{figure}

\begin{knowledgebox}{Lectura de la figura: conservar el lenguaje como entrada y quitarlo como cuello de botella intermedio}
La ruta de la izquierda traduce las señales de los sensores a lenguaje y luego decodifica la trayectoria a partir del lenguaje; la ruta de la derecha hace que Vision+Language entren directamente al modelo y que la Action salga directamente. No se trata de prescindir del lenguaje, sino de no usarlo como señal de supervisión intermedia obligatoria.
\end{knowledgebox}

\subsection{Por qué el lenguaje se convierte en un cuello de botella}

Esta subsección desglosa primero la expresión tajante «veneno». El problema del lenguaje no es que su semántica sea inútil, sino que, en cuanto se vuelve el intermediario obligado de las señales de conducción de bajo nivel, comprime el mundo continuo en etiquetas discretas.

El lenguaje es una abstracción humana de alto nivel, muy adecuada para instrucciones, explicaciones y sentido común; sin embargo, las señales de los sensores en la conducción autónoma contienen una gran cantidad de información continua, geométrica, dinámica y de grano fino. Si primero hay que traducirlas a lenguaje, se pierde mucha información continua y el proceso de entrenamiento se convierte en etiquetado humano o refinement humano. Para la conducción autónoma, que necesita data scaling masivo, esto limita la eficiencia en el uso de los datos.

\begin{importantbox}{El significado preciso de «Language es veneno»}
El lenguaje como entrada es muy valioso; el lenguaje como representación intermedia obligatoria puede ser tóxico. La toxicidad proviene de depender en exceso de la semántica humana, de reducir el aprovechamiento de los datos continuos, de formar un bottleneck y de obstaculizar el entrenamiento autosupervisado o con datos a gran escala.
\end{importantbox}

\subsection{«Quitar la L» y el aprendizaje autosupervisado}

Una de las lecciones importantes del éxito de la IA en el pasado es usar los datos para unsupervised/self-supervised learning. El aprendizaje autosupervisado construye señales de entrenamiento a partir de la propia estructura de los datos y reduce la dependencia del etiquetado humano. Si la conducción autónoma quiere seguir ese camino, debe reducir al mínimo las capas intermedias de supervisión humana y hacer que el modelo aprenda directamente de enormes volúmenes de datos visuales, de sensores y de trayectorias de conducción.

\begin{warningbox}{No hay que entender «quitar el lenguaje» como «prescindir de la capacidad lingüística»}
La conducción autónoma sigue necesitando comprender instrucciones de navegación, la semántica del tránsito y las intenciones humanas. Lo que se quita es el cuello de botella intermedio, no la entrada lingüística en la interacción humano-máquina ni la comprensión semántica dentro del modelo.
\end{warningbox}

\subsection{Resumen de la sección}

«Language es veneno» expresa una cautela frente a las representaciones intermedias. Para que la IA física aproveche plenamente los datos continuos de los sensores, el lenguaje no puede ser el paso estrecho por el que pase toda la información. Un VLA/Action decoding más directo busca mejorar la eficiencia en el uso de los datos y la capacidad de scaling.
\section{Eliminar la L y Software 3.0: no se trata de cambiar de nombre, sino de cambiar la ruta de los datos}

Las secciones anteriores trataron por separado la evolución de la pila de software y la eliminación del componente Language; esta sección examina ambos temas en conjunto. Liu Xianming (刘先明) recuerda en varias ocasiones que muchos términos nuevos hacen, en esencia, algo similar: modelos más grandes, más datos, menos reglas manuales y un paso más directo de la entrada a la acción. Lo que se denomina Software 3.0, VLA, VLM o de extremo a extremo no busca crear conceptos, sino encontrar rutas de datos que admitan más scaling.

\subsection{Por qué las estructuras intermedias reaparecen una y otra vez}

Esta sección explica una paradoja de ingeniería: las estructuras intermedias suelen agregarse para mejorar la controlabilidad, pero cuando el sistema alcanza cierto grado de complejidad, terminan limitando el techo del sistema de aprendizaje.

En los sistemas de conducción autónoma, las estructuras intermedias resultan atractivas: resultados de percepción, líneas de carril, cuadros delimitadores de objetos, trayectorias candidatas, go point, meta action, language token. Hacen que el sistema sea más interpretable, más fácil de depurar (debug) y más acorde con la división tradicional del trabajo de ingeniería. Sin embargo, cada estructura intermedia que se agrega puede introducir un cuello de botella de diseño manual: los datos se comprimen en algún formato definido por personas y el modelo solo puede aprender dentro de ese formato.

\begin{knowledgebox}{Ventajas y desventajas de las estructuras intermedias}
\begin{tabularx}{\textwidth}{p{0.22\textwidth}p{0.34\textwidth}X}
\toprule
Dimensión & Ventaja & Costo \\
\midrule
Interpretabilidad & Las personas entienden la salida de cada módulo, lo que facilita localizar problemas & Puede sacrificar información continua original. \\
División del trabajo & Percepción, predicción y planificación y control pueden avanzar en equipos separados & Las interfaces entre módulos se vuelven el techo del sistema. \\
Puesta en producción segura & Es más fácil aplicar reglas de protección y respaldo humano & Cuantas más reglas, más lento es el scaling de datos. \\
Aprovechamiento de datos & Las etiquetas intermedias permiten aprendizaje supervisado & Las etiquetas manuales limitan la escala y la capacidad expresiva. \\
\bottomrule
\end{tabularx}
\end{knowledgebox}

\subsection{El verdadero objetivo de eliminar la L: que los datos hablen por sí mismos}

A continuación se retoma la decisión de Xiaopeng (小鹏): eliminar la L no busca aparentar radicalidad, sino permitir que datos de mayor escala, más originales y más continuos entren en la cadena de entrenamiento.

El objetivo de eliminar el Language Bottleneck no es volver el sistema más misterioso, sino permitir que los sensores, las trayectorias de conducción y la retroalimentación de las acciones entren al entrenamiento con menos transformaciones manuales. El lenguaje puede usarse para instrucciones y semántica de alto nivel, pero si cada acción de bajo nivel tiene que pasar por language token, discretiza un mundo continuo y obstaculiza el aprendizaje autosupervisado y el aprovechamiento de datos a gran escala.

\begin{importantbox}{Definición mínima de Software 3.0}
En estas notas, Software 3.0 designa el uso de modelos grandes, grandes volúmenes de datos, señales autosupervisadas y fábricas de modelos en la nube para llevar las tareas del mundo físico de sistemas de reglas manuales a sistemas de aprendizaje con scaling sostenible.
\end{importantbox}

\subsection{Resumen de la sección}

Lo esencial de Software 3.0 no es la terminología, sino que la ruta de los datos sea más corta. Cuantas menos capas intermedias manuales haya, más posibilidades tiene el modelo de aprender directamente de datos reales; pero cuantas menos capas intermedias haya, más se necesitan evaluación, seguridad y gobernanza de ingeniería más sólidas.
\section{La transformación de XPeng (小鹏) hacia Physical AI: ciclo de datos del fabricante de automóviles y objetivo de Robotaxi}

Las secciones anteriores trataron la pila tecnológica y la eliminación de Language; esta sección examina por qué XPeng elevó la conducción autónoma al nivel de Physical AI. Liu Xianming (刘先明) afirma que «XPeng es, en esencia, una empresa de IA». Esta frase no es un eslogan publicitario; significa lo siguiente: el automóvil es uno de los agentes inteligentes del mundo físico de mayor escala, más controlables y con un ciclo comercial más completo; el fabricante de automóviles cuenta con usuarios reales, carreteras reales, retroalimentación real y una cadena de hardware controlable.

\begin{figure}[H]
\centering
\includegraphics[width=0.90\textwidth]{physical-ai-strategy.png}
\caption{La transformación de XPeng hacia Physical AI: de la capacidad de conducción autónoma a la inteligencia en el mundo físico. Diagrama conceptual propio, elaborado a partir de la conversación de 00:33:53--00:54:30.}
\end{figure}

\begin{knowledgebox}{Lectura de la figura: el automóvil es uno de los mayores cuerpos físicos de Physical AI}
En la figura, Physical AI no comprende solo el modelo, sino también el automóvil, los datos, el Robotaxi, la organización y la producción en serie. La ventaja de XPeng no está únicamente en los algoritmos, sino en que, como fabricante de automóviles, puede integrar los datos reales de su flota y la retroalimentación de la producción en serie en la iteración del modelo.
\end{knowledgebox}

\subsection{La ventaja en datos del fabricante de automóviles}

La mayor ventaja del fabricante de automóviles es una cadena de datos controlable. La flota circula en ciudades reales y genera datos de sensores, decisiones de conducción, retroalimentación de los usuarios, corner case anómalos, tomas de control y órdenes de incidencias. Siempre que el retorno de datos y la infraestructura de entrenamiento sean lo bastante buenos, el fabricante puede formar el ciclo «datos del mundo real -> entrenamiento en la nube -> despliegue en el vehículo -> retroalimentación real».

\begin{figure}[H]
\centering
\includegraphics[width=0.96\textwidth]{data-infra-loop-xpeng.png}
\caption{Ciclo de datos e Infra: la ventaja del fabricante de automóviles está en la cadena de datos controlable y en la retroalimentación de la producción en serie. Diagrama conceptual propio, elaborado a partir de la conversación de 00:02:16--00:19:00 y 00:33:53--00:54:30.}
\end{figure}

\begin{knowledgebox}{Lectura de la figura: el ciclo de datos no se reduce a «tener automóviles»}
La flota es solo el punto de entrada de los datos; un ciclo real requiere además la detección de corner case, infra de entrenamiento y análisis, actualización del modelo, despliegue en el vehículo y validación en producción en serie. Sin este sistema, los datos no se convierten por sí solos en capacidad.
\end{knowledgebox}

\subsection{El Robotaxi como hito de etapa}

Liu Xianming mencionó que espera que, en uno a tres años, el Robotaxi funcione muy bien en Guangzhou. El sentido de este objetivo es convertir el debate técnico en un servicio real: deja de ser una noticia del momento, un juguete o una atracción turística, y pasa a formar parte de lo que los usuarios usan por defecto cada día para desplazarse. La transición de Waymo en San Francisco, de atracción turística a servicio cotidiano, es la referencia que más le importa.

\begin{knowledgebox}{De atracción turística a servicio cotidiano}
Cuando Liu Xianming viajaba en Waymo con sus hijos en San Francisco, los niños lo llamaban Robot Car. Al principio, mucha gente iba a San Francisco solo para probarlo por novedad; pero cuando empezó a atender una gran cantidad de viajes diarios, y al salir de Caltrain Station ya había personas esperando un vehículo, dejó de ser solo una demostración tecnológica y pasó a formar parte de la vida de la ciudad. El objetivo de Robotaxi de XPeng también tiene que cruzar ese umbral de percepción.
\end{knowledgebox}

\subsection{Por qué Physical AI no es solo conducción autónoma}

Esta sección sitúa el objetivo de Robotaxi dentro del marco más amplio de Physical AI. La conducción autónoma es uno de los escenarios de la IA física más maduros, con más datos y con el ciclo comercial más claro, pero no es el destino final. El sistema del vehículo ya incluye percepción, predicción, planificación, control, interacción con el usuario, mapas, entrenamiento en la nube y mecanismos de seguridad; en el futuro, estos módulos se trasladarán a más agentes inteligentes físicos. Si XPeng logra demostrar el scaling en el cuerpo físico del automóvil, tendrá la oportunidad de extender la misma metodología a tareas más amplias del mundo físico.

\begin{importantbox}{El automóvil es el campo de entrenamiento de Physical AI}
El automóvil es un agente inteligente de alto valor, alta frecuencia y fuertes restricciones del mundo real. Permite que la IA no solo responda preguntas en una pantalla, sino que aprenda a actuar en medio de carreteras, tráfico y responsabilidades de seguridad.
\end{importantbox}

\subsection{Resumen de la sección}

Lo esencial de la transformación de XPeng hacia Physical AI es situar la conducción autónoma dentro de un ciclo de inteligencia del mundo real: los datos del fabricante de automóviles, el entrenamiento en la nube, el despliegue en el vehículo, la validación con Robotaxi y la ejecución organizacional conforman juntos la nueva estrategia.
\section{Lo simple es bello: retirar las capas intermedias y también los niveles jerárquicos de la organización}

Lo anterior trató la eliminación en el plano técnico; esta sección pasa a la eliminación en el plano organizacional. El «lo simple es bello» de Liu Xianming (刘先明) no es un lema estético, sino un principio de gestión de ingeniería: para que un sistema complejo itere con rapidez, es necesario reducir las capas intermedias innecesarias.

Esta sección pasa del «retirar» técnico a la «simplificación» organizacional. Liu Xianming dice que le gusta lo simple y que, al asumir el cargo, su decisión más importante fue simplificar los procesos, simplificar las etapas de investigación y desarrollo, fusionar las tareas duplicadas y bajar la prioridad de lo innecesario. No le gusta limitarse a escuchar informes; prefiere revisar directamente los problemas de primera línea, el código y los resultados de los experimentos.

\begin{figure}[H]
\centering
\includegraphics[width=0.96\textwidth]{simple-is-beautiful.png}
\caption{Lo simple es bello: retirar una y otra vez las densas capas intermedias para que los datos entrenen directamente las capacidades. Diagrama conceptual propio, elaborado a partir de la conversación de 00:00:00--00:02:16 y 00:54:30--01:48:46.}
\end{figure}

\begin{knowledgebox}{Lectura de la figura: la simplificación técnica y la simplificación organizacional son lo mismo}
Retirar la dependencia del LiDAR, retirar las reglas de planificación y control, retirar las capas intermedias y retirar Language tienen como fin reducir cuellos de botella; en la organización, reducir niveles jerárquicos, revisar directamente los problemas y fusionar el trabajo duplicado también buscan reducir los cuellos de botella de información.
\end{knowledgebox}

\subsection{Organización de ingeniería plana}

En escenarios donde la IA física se combina con la producción en serie, la velocidad de decisión es importante. Liu Xianming busca aplanar la jerarquía: que los engineer de primera línea puedan exponer los problemas directamente y que los lead de equipo también puedan escribir código, revisar experimentos y hacer deep dive. Dentro de la organización no se levantan muchos muros entre departamentos, para evitar el duplicate work, y la velocidad global aumenta al compartir experimentos y fracasos.

\begin{figure}[H]
\centering
\includegraphics[width=0.96\textwidth]{org-flat-engineering.png}
\caption{Organización de ingeniería plana: combinar la IA de frontera con la producción en serie exige decidir con rapidez sobre los problemas de primera línea. Diagrama conceptual propio, elaborado a partir de la conversación de 00:54:30--01:48:46.}
\end{figure}

\begin{knowledgebox}{Lectura de la figura: plano no significa sin gestión, sino cadena corta}
Los ingenieros de primera línea exponen código y experimentos, los responsables revisan directamente los problemas, el deep dive analiza los experimentos buenos y malos, la integración de recursos reduce la duplicación y, al final, todo ello permite un cambio rápido una vez que el modelo supera a las reglas.
\end{knowledgebox}

\subsection{Cultura de ingeniería: Deep Dive, compartir los fracasos y reducir la duplicación}

Esta subsección convierte lo «plano» en mecanismos cotidianos. Sin deep dive ni intercambio de fracasos, una organización plana corre el riesgo de tener solo unos cuantos niveles de reporte menos, en lugar de aumentar de verdad la velocidad de aprendizaje.

En la entrevista se menciona que el equipo hace deep dive: se comparten los experimentos buenos, los experimentos malos, lo que se ha visto recientemente y lo que a cada quien le salió mal. Esta práctica sirve a dos objetivos: primero, que el equipo forme rápidamente un contexto común; segundo, evitar que distintos grupos reinventen la rueda. La investigación y el desarrollo de Physical AI requieren una gran cantidad de experimentos; si los fracasos no se comparten, la organización cae una y otra vez en el mismo error.

\begin{importantbox}{Reutilización del conocimiento en la organización de ingeniería}
El activo central de un equipo de IA de frontera no son solo los pesos del modelo, sino también los experimentos fallidos, la experiencia en el ajuste de hiperparámetros, los métodos de procesamiento de datos y la comprensión compartida de los problemas. El Deep Dive vuelve explícito este conocimiento tácito.
\end{importantbox}

\subsection{Momento del cambio: oportunidades y riesgos}

A continuación se analiza el paso más arriesgado de la simplificación: cuándo puede eliminarse la ruta anterior. Para un fabricante de automóviles, esto no es un interruptor interno del grupo de investigación, sino una decisión que afecta directamente la seguridad de los usuarios y la calidad de la producción en serie.

Un fabricante de automóviles no puede lanzar algo solo porque la tecnología es atractiva. Para que el modelo sustituya de verdad a las reglas, hay que esperar a cierto punto: que el desempeño del modelo supere claramente al de las reglas, que los inconvenientes sean controlables, que los beneficios sean evidentes y que la responsabilidad de seguridad y la calidad de la producción en serie puedan sostenerse. Liu Xianming menciona el momento en que los probadores aplaudieron: esa es una señal organizacional de que el modelo superó a las reglas anteriores.

\begin{figure}[H]
\centering
\includegraphics[width=0.92\textwidth]{ai-transition-risk-balance.png}
\caption{Momento del cambio: oportunidades vs riesgos. El fabricante de automóviles debe esperar a que el modelo supere a las reglas, en lugar de lanzarlo solo porque la tecnología es atractiva.}
\end{figure}

\begin{warningbox}{El cambio de ruta tecnológica no puede imponerse por fe}
En un fabricante de automóviles, el cambio de ruta tecnológica debe responder por la seguridad, la calidad y los usuarios. Lo correcto no es «el modelo es atractivo, así que se lanza a toda la flota», sino lograr que el modelo supere al sistema anterior en los indicadores clave y ampliar el despliegue de forma gradual.
\end{warningbox}

\subsection{Resumen de la sección}

Esta transformación de Xiaopeng (小鹏) ocurre a la vez en dos planos, el técnico y el organizacional: en lo técnico se retiran las capas intermedias y en lo organizacional se acorta la cadena de información. Ambos sirven al mismo objetivo: que los datos y los problemas reales entren más rápido en la iteración del modelo.
\section{La perspectiva de la producción en serie: tecnología, calidad, costo y responsabilidad deben evaluarse en conjunto}

Antes se trataron la ruta tecnológica y la simplificación de la organización; esta sección agrega la perspectiva de la producción en serie. Liu Xianming (刘先明) subrayó varias veces que prefiere encontrar en los problemas reales lo que debe hacerse en la siguiente etapa. El responsable de conducción autónoma no es solo un responsable de investigación: también debe atender el producto, el negocio, la calidad, el hardware, el costo y la seguridad de los usuarios. Esta perspectiva explica por qué un fabricante de automóviles no puede copiar sin más el ritmo de un equipo de investigación.

\subsection{La producción en serie no es la etapa final de la investigación, sino una fuente de problemas}

Los equipos de investigación suelen definir primero la tarea y después buscar los datos y el modelo; el equipo de producción en serie enfrenta los problemas que los usuarios reales encuentran todos los días. Cada toma de control por parte del conductor, cada maniobra incómoda, cada error de juicio y cada queja de un usuario en la carretera se convierte en un problema que la siguiente iteración del modelo y de la ingeniería debe resolver. La producción en serie no es la fase de despliegue que sigue al término de la investigación, sino la entrada por la que surgen continuamente nuevos problemas de investigación.

\begin{knowledgebox}{Cómo los problemas de la producción en serie retroalimentan al modelo}
\begin{tabularx}{\textwidth}{p{0.22\textwidth}p{0.34\textwidth}X}
\toprule
Señal de la producción en serie & Qué hacer antes de que llegue al modelo & Valor de la retroalimentación \\
\midrule
Toma de control & Localizar la causa de la toma de control, extraer el fragmento y etiquetar el contexto & Encontrar los puntos ciegos del modelo y el límite de las reglas. \\
Incomodidad del usuario & Relacionarla con la maniobra, la velocidad, la distancia entre vehículos y el entorno & Mejorar la comodidad y el estilo de conducción. \\
Accidente/riesgo & Atribuir la responsabilidad y revisar la cadena de percepción/predicción/planeación & Formar muestras de seguridad y conjuntos de evaluación. \\
Restricciones de hardware & Registrar los límites de capacidad de cómputo, sensores, latencia y costo & Orientar la destilación, la cuantización y la poda. \\
\bottomrule
\end{tabularx}
\end{knowledgebox}

\subsection{Por qué «lo simple es bello» no significa eliminar sin criterio}

Simplificar no es eliminar todas las protecciones, sino quitar las estructuras repetidas, obsoletas o que obstaculizan el scaling. Un fabricante de automóviles todavía debe conservar la redundancia de seguridad, la validación de calidad y el proceso de liberación a producción. La simplicidad que subraya Liu Xianming consiste en que el proceso de desarrollo sea más corto, que haya menos objetivos y que los problemas salgan a la luz más rápido, no en sacrificar los límites de seguridad.

\begin{warningbox}{Simplificar no equivale a precipitarse}
Si, por perseguir el enfoque de extremo a extremo, se descuidan la seguridad, las pruebas y la calidad de la producción en serie, la simplificación se convierte en precipitación. La simplificación verdadera reduce la complejidad inútil y, al mismo tiempo, conserva las restricciones de seguridad necesarias.
\end{warningbox}

\subsection{Resumen de la sección}

La ruta de producción en serie de Physical AI exige que la tecnología y la responsabilidad avancen juntas. Los modelos deben ser más grandes y los datos más abundantes, pero cada liberación a producción debe superar la evaluación de seguridad, calidad, costo y experiencia de usuario.
\section{Detrás de los cambios de liderazgo: la misión histórica de cada etapa}

Las secciones anteriores trataron la tecnología y la organización; esta sección vuelve a los cambios de liderazgo. La conducción autónoma de XPeng (小鹏) ha tenido varios responsables: Gu Junli (谷俊丽), Wu Xinzhou (吴新宙), Li Liyun (李力耘) y Liu Xianming (刘先明). Liu Xianming no definió directamente la misión histórica de cada uno, pero puede leerse a partir de las etapas técnicas. En la etapa inicial había que explorar y construir capacidades. En la etapa intermedia se necesitaba un sistema basado en reglas y orientado a la producción en serie. Después vino la transición hacia el enfoque de extremo a extremo, y ahora se entra en la etapa de Physical AI y del scaling de modelos y datos más grandes.

\begin{figure}[H]
\centering
\includegraphics[width=0.96\textwidth]{leader-mission-stages.png}
\caption{Misión de los responsables de conducción inteligente de XPeng: cada etapa corresponde a una ruta técnica y a una tarea organizacional distintas. Diagrama conceptual propio, elaborado a partir de la conversación entre 00:54:30 y 01:48:46.}
\end{figure}

\begin{knowledgebox}{Lectura de la figura: el cambio de liderazgo no es un movimiento de personal aislado, sino un cambio de etapa}
Cada responsable enfrentó un problema principal distinto: la exploración inicial, la producción en serie basada en reglas, la transición hacia el enfoque de extremo a extremo, la Physical AI y, en la siguiente etapa, la integración en los servicios de la vida cotidiana. El cambio de liderazgo en una organización suele acompañar cambios en el stack técnico y en la estrategia de la empresa.
\end{knowledgebox}

\subsection{El papel de He Xiaopeng}

Liu Xianming describe a He Xiaopeng (何小鹏) como un jefe franco, apasionado por la tecnología y con preguntas muy sharp. Su primera conversación no fue una entrevista común, sino sobre «cómo dejar atrás a los competidores en la siguiente generación». En la planeación técnica, He Xiaopeng equilibra de forma dinámica las oportunidades y los riesgos; cuando el desempeño del modelo y la oportunidad se vuelven lo suficientemente claros, ajusta su juicio. Este proceso no depende de «convencer al jefe», sino de lograr resultados concretos, de modo que la emergencia de la tecnología cambie las decisiones.

\begin{importantbox}{La emergencia técnica cambia las decisiones}
En una ruta técnica con mucha incertidumbre, lo más convincente no es la discusión verbal, sino que el desempeño del modelo cruce de pronto cierto umbral. La organización debe dejar espacio suficiente para que la tecnología correcta se manifieste de forma natural en las métricas y en la experiencia.
\end{importantbox}

\subsection{Presupuesto, riesgo y las preguntas del CEO}

He Xiaopeng hace preguntas muy directas, por ejemplo «¿por qué se gasta tanto dinero?». Esta pregunta no busca solo recortar el presupuesto, sino obligar al equipo a explicar cómo la inversión en modelos más grandes, más datos y más infra se convierte en capacidades desplegables, experiencia de usuario y resultados comerciales. El scaling de la Physical AI es muy costoso; el papel del CEO no es solo aprobar el presupuesto, sino cuestionar de forma continua las oportunidades, los riesgos y las ventanas de tiempo.

\begin{warningbox}{El scaling no es gratuito}
Modelos más grandes y más datos aportan capacidades, pero también traen costos de entrenamiento, costos de despliegue en el vehículo, complejidad organizacional y riesgos de calidad. Una empresa automotriz debe evaluar la ruta técnica junto con las finanzas, la producción en serie y la seguridad.
\end{warningbox}

\subsection{Lecciones organizacionales de Meta y Google}

Liu Xianming menciona que el ritmo semestral y la cultura de reorg de Meta se prestan a la experimentación rápida propia de internet, pero en la era de la IA muchas cosas requieren construir sistemas desde cero, y una reestructuración grande cada seis meses no necesariamente favorece ese trabajo. Google también cambió tras el regreso de Larry Page y Sergey Brin. Esta comparación muestra que la cultura de la empresa influye en los proyectos de IA de largo plazo: si el ciclo de evaluación es demasiado corto, al equipo le resultará difícil construir sistemas de base que tardan más de un año en validarse.

\begin{knowledgebox}{La escala temporal de la organización en la era de la IA}
Los productos de internet pueden probarse con alta frecuencia y descartarse rápidamente; la Physical AI y la infraestructura de modelos grandes suelen requerir ciclos largos de validación. Si una organización solo acepta beneficios de ciclo corto, perjudica la construcción de capacidades de base.
\end{knowledgebox}

\subsection{De uno a cinco años}

Las previsiones de Liu Xianming a uno, tres y cinco años están bien escalonadas: dentro de un año, la inteligencia estará presente de forma más amplia en la vida cotidiana, probablemente todavía con predominio de la IA digital; dentro de tres años, la IA física podría entrar de verdad en la vida cotidiana; a cinco años es difícil prever, porque el mundo avanza demasiado rápido. Para las personas y las organizaciones, lo que se puede hacer es aprender de forma continua, reconocer la inercia de pensamiento y hacer bien lo que, con el conocimiento actual, resulta más correcto.

\subsection{Resumen de la sección}

Detrás de los cambios de liderazgo está la transición de etapa de la conducción inteligente de XPeng: de las reglas al enfoque de extremo a extremo, y de ahí a la Physical AI. La tarea de un líder no es solo dirigir al equipo, sino lograr que la ruta técnica correcta prospere mediante la ingeniería y la organización.
\section{Términos clave: índice de palabras clave de este episodio}

\begin{longtable}{p{0.24\textwidth}p{0.34\textwidth}p{0.33\textwidth}}
\toprule
Término & Explicación breve & Papel en este episodio \\
\midrule
Physical AI & IA en el mundo físico, capaz de percibir, decidir y ejecutar acciones & Marco central con el que Xiaopeng (小鹏) pasa de la conducción autónoma a una estrategia de IA. \\
Robotaxi & Servicio de taxis de conducción autónoma & Objetivo de etapa que Liu Xianming (刘先明) quiere que funcione bien en Guangzhou (广州). \\
Software 1.0 & Pila de software dominada por reglas, optimización y lógica escrita a mano & La ruta inicial de la conducción autónoma. \\
Software 2.0 & Pila de software basada en redes neuronales y guiada por datos & Concepto base de la conducción autónoma de extremo a extremo. \\
VLM & Vision-Language Model, modelo de visión y lenguaje & Entiende visión y lenguaje, pero no necesariamente emite una action de forma directa. \\
VLA & Vision-Language-Action, modelo de visión, lenguaje y acción & Dirección importante de la pila tecnológica de Physical AI de Xiaopeng. \\
Language Bottleneck & Cuello de botella intermedio que comprime la información continua de los sensores en token de lenguaje & Capa intermedia clave que Liu Xianming propone eliminar. \\
Self-supervised Learning & Aprendizaje autosupervisado: construye la señal de entrenamiento a partir de los propios datos & Permite aprovechar datos a gran escala y reduce la dependencia del etiquetado manual. \\
Scaling & Mejorar la capacidad con modelos más grandes, más datos y más cómputo & Liu Xianming lo considera la bet clave de Physical AI. \\
Fábrica de modelos en la nube & Entrenar modelos grandes en la nube y luego comprimirlos para desplegarlos en el vehículo & Estructura de ingeniería que conecta el scaling con la producción en serie en el vehículo. \\
Destilación & Usar un modelo grande para entrenar uno pequeño, de modo que el pequeño herede sus capacidades & Medio importante para el despliegue de la nube al vehículo. \\
Cuantización/poda & Reducir la precisión del modelo o recortar parámetros para adaptarlo al hardware & Necesaria para el despliegue en el vehículo y el control de costos. \\
Infra & Infraestructura de datos, entrenamiento, análisis y despliegue & Base que destacan tanto la ruta de Cruise como la de Xiaopeng. \\
Corner Case & Escenarios límite de cola larga & Objeto que el ciclo cerrado de datos de la conducción autónoma debe seguir buscando. \\
\bottomrule
\end{longtable}

\subsection{Resumen de la sección}

Los términos de este episodio giran en torno a una pregunta: cómo llevar la conducción autónoma de la ingeniería basada en reglas a un sistema de Physical AI capaz de hacer scaling. Las palabras clave no son sustantivos aislados, sino una cadena que une entrenamiento, compresión, despliegue y retroalimentación de la producción en serie.
\section{Síntesis y ampliación}

\subsection{Conclusiones centrales}

\begin{enumerate}
\item El giro de XPeng (小鹏) hacia Physical AI no consiste en añadirle un modelo grande al automóvil, sino en reescribir el ciclo cerrado de datos y modelos de la conducción autónoma.
\item La trayectoria de Liu Xianming (刘先明) muestra que los problemas del mundo real y el sentido de misión orientan de manera constante las decisiones técnicas.
\item La pila de software de conducción autónoma ha pasado de las reglas, a los sistemas parcialmente basados en modelos y a los de extremo a extremo, y avanza hacia modelos más grandes y fábricas en la nube.
\item El significado preciso de «Language es veneno» es que el lenguaje, como cuello de botella de supervisión intermedia, limita el data scaling.
\item El objetivo de eliminar la capa intermedia es que los datos continuos de los sensores y las salidas de acción entren de forma más directa en el entrenamiento.
\item La ventaja de los fabricantes de automóviles está en sus flotas reales, en una cadena de datos controlable y en la retroalimentación de la producción en serie.
\item La fábrica de modelos en la nube se encarga de entrenar los modelos grandes, que después se destilan, cuantizan y podan para desplegarse en el vehículo.
\item La simplificación técnica y el aplanamiento de la organización son lo mismo: ambos reducen cuellos de botella.
\item Detrás del cambio de liderazgo hay un cambio de etapa técnica; la tarea de Liu Xianming es convertir la ruta de Physical AI en un servicio para la vida cotidiana.
\item La bet clave para el futuro es el scaling en Physical AI, junto con identificar de forma continua los errores del propio pensamiento.
\end{enumerate}

\subsection{Preguntas abiertas}

Las preguntas abiertas se dejan al final porque la ruta de XPeng sigue desplegándose a gran velocidad. Lo que realmente hay que observar no es solo «qué se eliminó», sino también si, después de eliminarlo, el sistema es más estable, más controlable y más capaz de escalar.

\begin{itemize}
\item Tras eliminar Language, ¿cómo conserva el modelo suficiente semántica de alto nivel e interpretabilidad?
\item ¿La destilación del modelo grande en la nube al modelo del vehículo se convertirá en el límite superior de la experiencia?
\item ¿Puede el ciclo cerrado de datos de un fabricante de automóviles alcanzar el ciclo cerrado de alta calidad de las flotas dedicadas de Robotaxi?
\item Si Robotaxi funciona bien en Guangzhou, ¿qué obstáculos no técnicos quedan para escalarlo a todo el país?
\item ¿Cuánto tiempo puede sostenerse una organización de ingeniería plana dentro de la producción en serie a gran escala y de un sistema de calidad?
\item ¿El scaling de Physical AI se validará primero en el automóvil o en robots y otros cuerpos físicos?
\end{itemize}

\subsection{Lecturas adicionales}

\begin{itemize}
\item Entrevista de EP121 con Tan Jie (谭杰): robots, modelos del mundo, transferencia entre cuerpos físicos y Gemini Robotics 1.5.
\item Entrevista de EP132 con Gao Jiyang (高继扬): Waymo, Momenta y la llevada a producción de la inteligencia corporeizada en Galaxea (星海图).
\item Panorama de datos de EP134: Data Recipe, datos para robots y fijación de precios de los datos.
\item Materiales sobre Software 2.0, VLA, VLM, aprendizaje autosupervisado y destilación de modelos.
\end{itemize}

\begin{importantbox}{Valoración final}
Lo que más vale conservar de EP120 no es la frase «quitar la L», sino una filosofía de ingeniería: cuando una capa intermedia se vuelve un cuello de botella para el data scaling, hay que eliminarla; cuando los niveles jerárquicos de la organización se vuelven un cuello de botella para la retroalimentación de los problemas, también hay que eliminarlos. El avance de Physical AI proviene de datos más directos, modelos más grandes, cadenas de ingeniería más cortas y una validación más decidida en la producción en serie.
\end{importantbox}

\end{document}
